%==============================%
%     NORMAL TEXT VARIABLE     %
%==============================%

% --- Judul ---
%
% Gunakan \judul di bawah ini untuk mengganti judul bawaan:
%  - Lembar Jawaban <SKEMA TUTORIAL>
%  - Tugas <SKEMA TUTORIAL> <NOMOR TUGAS>
% Menjadi judul apa pun yang Anda inginkan
%
%\newcommand{\judul}{Judul Tulisan Anda}

% --- Skema Pembelajaran/Tutorial ---
% [Syntax] \newcommand{\skemaTutorial}{PERINTAH atau TEKS}
%
% Isilah PERINTAH dengan yang ada di bawah ini untuk penerjemahan
% saat bahasa diganti menjadi Bahasa Indonesia atau Bahasa Inggris
%
%  - \TutorialOnline
%    Tutorial Online <--> Online Tutorial
%
%  - \TutorialWebinar
%    Tutorial Webinar <--> Webinar Tutorial
%
%  - \TutorialTatapMuka
%    Tutorial Tatap Muka <--> Face-to-Face Tutorial
%
% Anda bisa saja mengganti dengan teks Anda sendiri, namun
% penerjemahan nama Skema Tutorial tidak akan bekerja
%
\newcommand{\skemaTutorial}{\TutorialOnline}

% --- Informasi Mata Kuliah, Sesi, Tugas, dan Tutor ---
\newcommand{\sesiKe}{3}
\newcommand{\tugasKe}{1}
\newcommand{\tanggalLengkap}{\today}
\newcommand{\tahun}{\the\year}

\newcommand{\namaMataKuliah}{Pengendalian Berbasis Desktop}
\newcommand{\kodeMataKuliah}{STSI9999}
\newcommand{\KodeMataKuliah}{STSI-9999}

\newcommand{\kodeKelas}{256}

\newcommand{\namaTutorPengampu}{Revo Wibowo, S.Pd., M.Pd., M.Sc.}

% --- Informasi Mahasiswa ---
\newcommand{\namaMahasiswa}{Yoeru Sandaru}
\newcommand{\nimMahasiswa}{081298765432}
\newcommand{\programStudi}{Sains Data}
\newcommand{\kodeProgramStudi}{/257}
\newcommand{\fakultas}{Sains dan Teknologi}
\newcommand{\kodeFakultas}{/127}
\newcommand{\utDaerah}{UT Medan}
\newcommand{\kodeUTDaerah}{/28}
\newcommand{\perguruanTinggi}{Universitas Terbuka}